\section[Objective 1: Develop SDR-to-SBC Interface for Radar Control and Data Acquisition]{Objective 1: Develop \gls{sdr}-to-\gls{sbc} Interface for Radar Control and Data Acquisition}

The methodology for interfacing the \gls{sdr} with the \gls{sbc} follows a structured software engineering lifecycle. It includes evaluation of the legacy interface, software implementation, unit testing, and system-level validation. All project source code will be authored, reviewed, tested, and version-controlled by the project team.

\subsection{Research and Evaluation of Existing Interfaces}

The first phase reviews the software interfaces provided by the radar manufacturer \cite{luswavePUPDUAL24C}. This includes analysing the legacy MATLAB-based drivers to understand the communication logic and determine whether existing libraries can be reused. The outcome will determine whether a custom \gls{api} must be developed from scratch or whether the legacy drivers can be adapted for the target hardware. These findings will also define the technology stack for implementation.

\subsection{Software Implementation}

Following the research phase, an API will be developed and implemented. Development is divided into two stages: establishing the communication protocol and implementing the radar command structure. First, the \gls{usb} firmware protocol will be implemented to manage the data link between the radar and the \gls{sbc}. The radar control commands, which configure transmit and receive parameters, will then be mapped and integrated into the API. To support deployment on the target platform, the API will be built as a static library compatible with Ubuntu 22.04. Due to the large variety of radar commands involved and the limited project timeframe, Claude Code will be used to assist in implementing the different commands once the command protocol structure has been designed. A human-in-the-loop approach will be adopted, where all outputs generated by Claude Code are reviewed and verified to ensure correctness.

At the end of this stage, the \gls{usb} communication protocol will undergo preliminary testing on a standard Linux computer connected to the radar. The radar control commands will also receive basic functional verification. This testing will remove compile-time errors and perform initial validation checks, such as confirming that captured data volume is consistent with the configured sampling rate. Rigorous data-integrity checks and command-level validation will be completed during unit testing.

\subsection{Unit Testing}
\label{chap:413}

Initial unit testing will be conducted in a standard Linux desktop environment to isolate variables associated with embedded hardware. The radar receiver channel will be connected to a known signal generated by a signal generator, while the transmitter channel will be connected to a frequency analyser to observe the transmitted output. Individual radar configuration parameters, such as sweep time, sampling frequency, and active channels, will be varied one at a time. Each command will then be validated by inspecting the API output and laboratory equipment measurements to verify the transmitted and received radar waveforms, ensuring that the radar parameters update correctly.

\subsection{Ground Integration and Testing}
\label{chap:414}

Once the API passes desktop unit testing, the static library and associated executables will be ported to the target \gls{sbc}, a Raspberry Pi. The \gls{sbc}, software interface, and power system will then be integrated onto a moving ground platform. Ground testing will include a controlled motion Doppler test, where the radar is mounted on a ground vehicle moving at a constant velocity. In this scenario, the measured Doppler shift will be compared against the expected velocity. Testing will verify that valid radar data is generated and stored in the designated Raspberry Pi memory locations under operational scenarios.

Additionally, stress-test scenarios, such as extended-duration data capture, will be conducted to evaluate system stability. Data throughput will also be monitored using Linux command-line tools, such as iperf3, to ensure the interface can reliably process all incoming data without packet loss or dropped samples \cite{iperf_fr}.

Objective 1 will be complete once the system passes the ground integration test. The transition from ground testing to flight testing is not expected to change the software interface, as it operates as a self-contained component within the wider system.